\chapter{A Curriculum, Read Cautiously}
\label{ch:curriculum}

\begin{plainlanguage}
\textbf{In plain terms:} The Kalabari life cycle looks, at first, like a tidy curriculum: easy lessons for children, hard willpower-testing ones for adults. A careful look at the early examples breaks that tidy story — willpower-testing lessons show up surprisingly early. The honest reading is a curriculum that blends both kinds of lesson from the start.
\end{plainlanguage}

\section{A pattern across the life cycle}

Read across Wariboko's five life-cycle stages --- \emph{awome} (childhood), \emph{asawo} (young adulthood), \emph{opuasawo} (elders), \emph{alapu} (chiefs) and their \emph{iya} wives, and finally the status of ancestor --- the examples suggest a developmental pattern. Early counterfoils are overwhelmingly high-$d_p$, low-$m$: the wrongness is visually self-evident and not particularly tempting. Counterfoils involving the apprentice trader's restricted commercial freedom, or the elder's restraint in administering punishment, occupy a middle zone where the boundary requires understanding of social and institutional structure rather than physical perception, and where real stakes begin to apply pressure. By the time one reaches chieftaincy installation, bibife, and the conduct expected of a dying person facing their final moments, $d_p$ has fallen close to zero --- no physical absurdity remains to lean on --- while $m$ remains high or increases further.

\begin{definition}[Curriculum Trajectory]
A stage-indexed map $s \mapsto (d_p(s), m(s))$ tracking how perceptual margin and motivational pull change across an individual's developmental stages.
\end{definition}

\begin{naivemodel}
It is tempting to state this as a clean, monotonic curriculum: $d_p$ falling and $m$ rising in lockstep across five neat stages, with recognition training handed off cleanly to implementation training partway through. The temptation should be resisted, and resisting it here is an instance of exactly the discipline Part~III's two case studies were built to model: a dozen well-chosen ethnographic anecdotes, selected by Wariboko to illustrate a different thesis (the aristocratic ideal), can be made to fit more than one curriculum shape, and treating any one shape as established by the examples rather than merely consistent with them would repeat the error both Chapter~\ref{ch:room-for-river} and Chapter~\ref{ch:development} were built to diagnose and correct.
\end{naivemodel}

Closer inspection of the \emph{awome}-stage material itself undermines strict sequentiality: the drink-sharing and shared-plate rituals are pure implementation challenges (low $d_p$, real $m$) embedded early, well before a clean handoff point would predict. The more defensible reading is interleaved rather than sequential.

\begin{repairbox}
\begin{conjecture}[Interleaved Exposure]
Rather than a strict handoff $I_R \to I_I$, effective training exposure at stage $s$ is better modeled as a weighted combination,
\[
E_s = \alpha_s I_R + \beta_s I_I,
\]
with $\alpha_s > \beta_s$ early in development, and the ratio shifting toward $\beta_s$ as the individual matures.
\end{conjecture}

This is stated as a conjecture, not a theorem, and not even a fully established proposition: it is motivated by, but not established by, the same dozen non-randomly-sampled illustrative cases just discussed. The anti-vacuity principle introduced in Chapter~\ref{ch:development} applies here directly and is worth restating explicitly rather than left implicit: this curriculum hypothesis is admissible to entertain because the recognition/implementation distinction it relies on was independently motivated in Chapter~\ref{ch:two-axes}, before this curriculum question was asked --- not invented here to make Wariboko's examples fit a pre-decided shape.
\end{repairbox}
